\begin{abstract}
Этот том объединяет пять статей исследовательской программы
Songlines в один документ. Программа спрашивает, как агенты с
символической семантической памятью выполняют навигационные задачи
--- и что меняется, когда память становится коллективной.
\textbf{Часть I} вводит стадийную декомпозицию Q/R/M/C (Query $\to$
Retrieval $\to$ Materialization $\to$ Completion): один
logging-контракт, локализующий отказы, которых не видят метрики
успеха, переносимый на пять классов систем --- от символических
одиночных агентов до LLM-стеков; острейший одноагентный результат ---
связывающее ограничение не в ранжировании, а в \emph{генерации}
кандидатов ($91\%$ пустых извлечений при точности ранжирования
$96{,}9\%$). \textbf{Часть II} делает «навигацию по чужой памяти»
измеряемым событием: каждое обязательство аннотировано чужой долей
$\varphi$ его свидетельств; сырые чужие локи почти никогда не
завершаются напрямую, но окупаются как транспорт; закон дистанции в
замкнутой форме предсказывает горизонты намерения с точностью до
тика; а двухстрочный координационный контракт (Warp Drive) делает
самый быстрый режим обмена и лучшим. \textbf{Часть III} находит
переносимую единицу: не точку (координату), а песню --- маршруты с
провенансом рёбер, наследующие завершение и безопасность свидетеля и
рвущиеся на предсказанном ребре; идентичность мест, восстановленную
из одного смысла (отказывающую закрыто там, где координаты отказывают
открыто); и результат анатомии: узлы несут идентичность, рёбра ---
перенос, за $6{,}4\%$ бит снапшота. \textbf{Часть IV} даёт
коллективной памяти механику --- слияние, доверие и устаревание как
три явных правила, занимающие новый режим на плоскости $M{\times}C$
--- и математику: категорную рамку, в которой морфизмы выравнивания
конструируются, а не предполагаются, а фазовая диаграмма
доверие$\times$каденция решает, охватывает ли копредел коллектив.
\textbf{Часть V} обращается к \emph{онтогенезу} памяти: формирование
как двухосевое решение (контрфактическая маргинальная полезность
$\times$ простота аналогии) с операцией исключения, защищающей от
ложной генерализации, --- измерено детерминированно, выучено
контекстным бандитом, обыгравшим собственного проектировщика, и
эволюционировано к конечному бюджету куплетов. Затем Часть V собирает
все механизмы в один рантайм и закрывает acceptance-тест рецензента:
runtime-оцениватель полезности, совпадающий с точным контрфактом
(Спирмен $0{,}989$), так что контроллер не нуждается в оракуле;
допуск с валидацией на стороне приёмника, делающий социальный обмен
нетто-полезным на горизонте, где нерегулируемый обмен вреден;
сравнение при равных бюджетах, где сырая история не догоняет и при
$4\times$ памяти; тест на adversarial-провенанс; safety-гейты,
доводящие fail-open ниже $1\%$ на сеточном и континуальном
субстратах; и полный учёт стоимости провода. Части писались как
самостоятельные подачи; перекрёстные ссылки между ними появляются
здесь как части одного тома.
\end{abstract}

\section*{Как читать этот том}

\paragraph{По одному предложению на часть.} Часть I строит инструмент
(стадийное измерение); часть II им расценивает чужие свидетельства
(провенанс и его закон); часть III находит правильную единицу
переноса (песня: ориентиры $+$ такты); часть IV даёт коллективную
механику и математику «иметь в виду одно и то же»; часть V
спрашивает, откуда вообще берётся память, которую координируют части
I--IV, и делает правило её формирования выучиваемым и
эволюционируемым.

\paragraph{Происхождение.} Части I и IV выросли из одной ранней
рукописи (\emph{Q/R/M/C Framework $+$ Minimal Collective Semantic
Memory}, хранится в архиве серии как \texttt{symbolic\_memory});
каждый эксперимент предка присутствует здесь, в уточнённом виде, в
частях I и IV. Части II и III --- компаньоны, построенные на
инструментарии части I. Там, где текст части говорит
«сопутствующая статья», имеется в виду другая часть этого тома.

\paragraph{Измерительная дисциплина.} Каждое утверждение каждой части
следует одному протоколу: предсказания регистрируются на диск до
исполнения, отказы относительно регистрации публикуются с
механизмами, а сильнейшие результаты --- точные совпадения с
предсказаниями в замкнутой форме ($6/6$ hold-out закона дистанции,
$110/110$ закона разрыва, $80/80$ восстановления $SE(2)$, $0$
fail-open в $252$ ячейках абляции словаря).

\paragraph{Код и данные.} Все эксперименты живут в одном репозитории
(\texttt{experiments/}, \texttt{songline\_drive/}); приложение о
воспроизводимости каждой части перечисляет точные команды. Сырые
артефакты прогонов --- в \texttt{tmp/} с файлами \texttt{RESULTS.md}
у ключевых экспериментов.
